% Candidate appendix methods, not included in the active manuscript.
\subsection{Measurement Availability and Experimental Scope}

In the analysis figures, state $t$ means $t$ completed parameter updates.
Loss $L_t$ and raw gradient $g_t$ are evaluated at $\theta_t$; $m_t$ denotes the
stored first moment after consuming this gradient, available before the next
parameter displacement. In the manuscript's update-index convention, these
vectors are $g_{s+1}$ and $m_{s+1}$ at parameter state $s$. An integration must
use one notation consistently; the numerical availability clock is identical.
Losses cover states 0--1,024; stored moment diagnostics cover 0--1,023.
No next-step loss, gradient or parameter increment is supplied to a forecast.

The MLP uses the frozen 5,000-image CIFAR-10 objective; ResNet-32 uses all 50,000
training images with the 50 fixed ghost-batch-normalization groups of 1,000.
The Transformer uses all 2,088,600 prepared WikiText-2 target tokens per update.
The frozen model/data/cache identities and exact initialization hashes are
retained in Run 215. Rates are $10^{-4},3\cdot 10^{-4},10^{-3}$ for the MLP and
$3\cdot 10^{-4},10^{-3},3\cdot 10^{-3}$ for the other cases. Training is FP32
with FP64 diagnostic reductions and TF32 disabled. No Hessians are measured.
All quantities concern the full specified training objective and all trainable
parameters, counting tied parameters once. Zero-norm cosines/projected changes
are undefined rather than regularized; the retained forecast rows are finite.

\subsection{Loss-Only Episodes and Warning Windows}

Searching begins at state 64. At each unassigned state, the episode reference
is the minimum over the preceding 16 states, excluding the current state.
An episode starts when current loss is at least 1.10 times that reference.
Freeze the reference until the first later loss strictly below it, merge all
intermediate rebounds, and resume searching after recovery. The largest loss
before recovery defines the observed peak. If no recovery is observed by 1,024,
retain the episode and mark its recovery as right-censored. A large episode
has an observed peak at least 1.50 times the frozen reference. This is an
observed-horizon classification, not a claim about an unobserved final peak.

An early warning must be strictly before onset, in the preceding 8 states,
with loss less than 1.05 times the minimum over its own current and previous 16
states. The alarms are projected change above 2; a fresh crossing from at most 2
to above 2; one-state gradient-norm growth above 20\%; one-state stored-momentum
growth above 20\%; and the conjunction of fresh crossing and momentum growth.
Record hits, misses and the earliest qualifying lead within the window.
Lead summaries condition on a hit; they do not assign zero to misses.

The micro examples select each highest-rate arm's completed episode with the
largest peak/reference ratio, using loss alone. Crops begin 10 states before
the peak and end at the first subsequent loss below the crop-start value.
That display reference differs from the episode's frozen minimum. The repeated
event view retains all 27 large highest-rate episodes, with no diagnostic-based
selection and no censoring exclusion. Every displayed relative state from-8
through+8 has data. Five MLP windows include the end of a previous episode;
they are retained and flagged. Each individual trace remains visible. The
pointwise median is an aggregate across dependent episodes, not a smoothed
trajectory or an estimate with an independent-event confidence band.

\subsection{Forward Evaluation and Controls}

At eligible state$t$, predict whether
\begin{equation}
  \max_{1\leq h\leq 8} L_{t+h}/L_t-1\geq 0.10.
\end{equation}
Eligibility requires$t\geq 64$ and current loss less than 1.05 times its minimum
over$t-16,\ldots,t$. Omit the final 8 states with incomplete future coverage;
do not label them negative. This target uses current loss and a fixed horizon,
whereas the episode census uses a frozen pre-onset baseline and variable recovery.
Adjacent decision targets overlap and are not independent observations.

The baseline features are nine arm indicators, $t/1,024$, log current loss,
log-loss changes over 1,4 and 8 updates, and current log draw-up above the
past/current 16-lag minimum. The primary extension adds continuous projected
change. Other inherited comparisons add log gradient level and one-state
growth, log first-moment level and growth, or all five diagnostic features.
All fits use training-only standardization and L2 logistic regression with
$C=1$, liblinear, fixed solver randomness, no class weighting and no tuning.

Test blocks begin 384,512,640,768 and 896, ending before the next block or at 1,016.
Training decision states are strictly below the test start minus 24. Hence all
training outcomes precede the earliest loss-history input in the test block.
The same time cuts apply across all arms; individual states are never randomly
split. Later blocks may train on previous test periods, as in expanding-window
forecasting, but never on their own future. Report pooled, per-case, per-arm
and per-block results. Per-case results subset predictions of the pooled fits.

Average precision is the noninterpolated precision-recall summary; Brier score
is mean squared probability error. Also report prevalence, decision counts,
positive counts, fixed-alarm precision/recall and false-positive rates.
An additional reference predicts the positive prevalence of its training set,
constant within each test block. It provides a probability-error benchmark,
not a probability calibration fit. No test operating threshold is selected.

The new exploratory checks add raw cosine alone and raw cosine with one-state
log gradient growth; replace the three selected loss lags with all 16 lagged
log-loss changes; and compare current projected change with its 16- and 32-state
delays. All delay comparisons, including their controls, use a common 48-state
gap so every training outcome precedes the earliest delayed test input.
Delaying a feature is not an exchangeable permutation null. These comparisons
provide descriptive timing sensitivity without a formal causal or significance
test. The inherited sensitivity grid includes all combinations of horizons
4/8/16 and future rises 5/10/25\%, with no best-setting selection.

\subsection{Scope of Inference}

This is a retrospective exploratory study motivated by an inspected peak.
The time blocks are held out from coefficient fitting, not pristine confirmation
of hypotheses chosen after inspecting the results. Repeated events, overlapping
forecast windows, rates sharing an initialization and successive training blocks
are dependent. We report descriptive counts, ranges and block variation, without
IID confidence intervals or confirmatory$p$-values. One initialization per case
does not establish robustness across initializations; full-batch training does
not remove that limitation. Initial-gradient qualification and common frozen
inputs do not isolate all later cross-hardware numerical effects. No inference
of Hessian curvature, universal Adam stability, generalization, causal mechanism
or optimizer-intervention benefit is made.
